\chapter{Background}
\label{sec:background}

This chapter introduces beni at the depth the rest of the thesis needs: the language, its effects
and its boundary with JavaScript, how its compiler is organised, how it represents lists today, its
release optimiser, and its two browser back ends. It ends with a prototype study of what in-place
writes would buy and a precise statement of the question.

\section{beni in brief}
\label{sec:bg-beni}

beni is a statically typed, strict, pure functional language in the family of Elm. Its compiler,
written in Zig, emits modern JavaScript as ES modules, and its primary platform is the browser; Node
is supported as well. It keeps Elm's guarantees---no run-time exceptions in ordinary code,
exhaustive pattern matching, managed effects---and departs from Elm in a few deliberate places.
Those that matter here are the following.

\paragraph{Syntax.} beni uses Unicode notation: \code{→} for arrows, \code{λ} for lambdas,
\code{▷} for the forward pipe, \code{…} for a spread, \code{⊤} for the unit type and value, and
\code{×} for tuple types. A function body is a block of bindings followed by its result. Lists are
written, built and matched with brackets and spreads: \code{[ x, …rest ]} matches a non-empty
list and binds its tail, \code{[ …init, last ]} matches from the other end, and \code{[ …xs, …ys ]}
concatenates. Records are immutable and updated by \code{\{ r | f = v \}}.
A small \code{update} function in the style of a TodoMVC implementation looks like this:
\begin{Verbatim}
update : Msg, Model → Model
update msg model =
    case msg of
        Typed text →
            { model | draft = text }
        Add title →
            { model | todos = model.todos ++ [ { id = model.nextId, title = title, completed = False } ]
                    , nextId = model.nextId + 1 }
\end{Verbatim}

\paragraph{No automatic currying.} Function types are $n$-ary---\code{Msg, Model → Model} takes two
arguments---and every call is saturated. Partial application is written with a placeholder,
\code{f a \_}, which is sugar for the lambda \code{λp → f a p}; the pipe \code{x ▷ f a} is rewritten
to \code{f x a} before anything runs; a bind \code{let x ← e} in a block is sugar for a written
lambda as well. There are no point-free composition operators. Every closure in a beni program is
therefore a lambda the programmer wrote (or a placeholder that desugars to one), and the compiler can
see what it captures. There is also no calling convention for partial applications in the emitted
code: every emitted JavaScript call passes all its arguments.

\paragraph{Evaluation order.} beni is strict and evaluates left to right in source order: in an
application, the callee and then the arguments, left to right, and then the call; record literals
field by field in written order; a record update evaluates the record and then the updated fields;
block bindings in the order written; a self tail call evaluates all its new arguments before any
parameter is rebound. The order is part of the language definition and is observable today through
\code{Debug.log}; the compiler's release optimiser is constrained not to change it
(\S\ref{sec:bg-release}).

\paragraph{Values and identity.} A record update \code{\{ r | f = v \}} builds a new record whose
\emph{other} fields are the very same values as \code{r}'s, not copies of them. The language offers
programs no way to compare two values by reference: \code{==} is structural. The runtime, however,
may compare by reference to skip work, and the browser runtimes do (\S\ref{sec:bg-browser}).

\paragraph{Static dispatch.} A type's methods are the public values of the module that declares
it; \code{x.m a} calls one, and a top-level annotation may carry a constraint such as
\code{where a.compare : a, a → Order}. The operators \code{==} and \code{<} call the receiver type's
\code{eq} and \code{compare}, derived structurally when the type declares none. Evidence for a
constraint is passed as a hidden argument. At a call site the checker resolves each method call to a
term of a dispatch table: a top-level function, a derived \code{eq} or \code{compare}, or evidence the
enclosing function received as a parameter. A top-level value that carries a \code{where} clause is
memoised once per evidence. The constraints of an unannotated declaration are inferred and published
in its module's interface.

\section{Effects, fibers and commands}
\label{sec:bg-effects}

\paragraph{Inferred effects.} beni has no effect syntax. Its type checker infers, for every function
type, two bits---whether a call may perform an effect, and whether it may suspend the current
fiber---and solves them as a directed graph of inclusion edges (``this class is at least as
effectful as that one'') after each module is type-checked. The bits live beside the unifier, not in
it: nothing is added to the type store or to unification. The solved bits are published in the
module's interface and hashed with it. A function that takes a callback is summarised
\emph{polymorphically} in the callback's bits: a reference to a top-level declaration instantiates its
summary, each class of the summary receiving a fresh class at the use that is related to the
summary's classes by inclusion edges. That is what lets \code{List.map} serve a pure and a suspending
callback with one definition. A function that may suspend is compiled into a suspendable form whose
fast path is an ordinary call and which hands a continuation to the runtime when it parks.

\paragraph{Fibers.} Concurrency is cooperative: fibers interleave only at suspension points, and
continuations are one-shot. The core \code{Task} module provides \code{spawn} (start a fiber),
\code{join} (wait for a fiber's outcome; every joiner receives the same value), \code{scope} (a
structured-concurrency block whose body runs at once and whose children are closed when it ends) and
\code{bracket} (acquire a resource, use it, and release it on every exit, the release being stored on
the fiber's list of finalisers). Fibers communicate through mutable cells (\code{Ref}, with
\code{get}, \code{set} and \code{update}) and one-shot deferred values (\code{Deferred}); the core
library has no queue today, and one added later would be treated as a \code{Ref} is. A cancelled
fiber never resumes the rest of its body.

\paragraph{Commands and subscriptions.} A browser program does not perform effects directly from
\code{update}: it returns commands, which are thunks the runtime stores and runs later
(\code{Cmd.perform (λsend → ⋯)} runs a function that may send messages back, and \code{Cmd.do} runs a
function for what it does). A command that may wait runs in a fiber, concurrently with later
messages; one that never waits is queued to run soon after \code{update} returns, and nothing in
the runtime's contract today orders it before the next message. Subscriptions are computed from the
model once per render and diffed by the runtime: a subscription whose key stayed keeps its fiber and
takes the new functions that turn its events into messages.

\paragraph{Defects.} A failure that the program cannot handle---a bug in the program or the host---is
a \emph{defect}. It stops the program: a development build shows a crash screen, and a release build
logs and stops. Cleanup registered by brackets still runs.

\paragraph{Debugging aids.} \code{Debug.log tag x} prints \code{x} and returns it; \code{Debug.toString}
renders a value as text. A release build refuses a program that reaches \code{Debug}, so that a release
build cannot behave differently from a development build through them.

\section{The boundary with JavaScript}
\label{sec:bg-boundary}

Only privileged platform packages and the embedded core library may declare a \code{foreign} value.
A \code{foreign} binds to an export of a \emph{sibling} JavaScript file of the same module, one export
per declaration, and the compiler checks the sibling's shape at build time: the type of a foreign
value has one of two permitted forms, every declaration has an export and every export a declaration,
imports are covered, and the export's parameter list has the declared arity. Every \code{foreign}
already declares its \emph{effect rung}---pure, impure or suspending---and functions passed to
platform code can be marked \code{sync}, meaning they may not suspend. The compiler trusts these
declarations: a wrong rung is a bug in privileged code, not a user error.

Below \code{foreign}, the core library has a low-level \code{Js} interface through which much of it
is written in beni itself: \code{Js.pure} wraps a computation on raw host values, \code{Js.call}
calls a method on a host value, \code{Js.from} turns a beni value into a host value, and \code{Js.to}
reads a host value back. Since most of the list library is written this way, the list writers'
bodies consist of \code{Js} calls that slice and write arrays.

One existing promise about lists at the boundary matters here: a sibling never writes a list it was
given, and never keeps one it will later write. It may keep a list it was given in order to read it
later, and some platform code does.

\section{Modules, interfaces and incremental builds}
\label{sec:bg-modules}

A beni program is a set of modules whose import graph is acyclic. Each module is type-checked once
its imports have been, and independent modules are checked in parallel. The result of checking a
module is an \emph{interface}: a record of the types of its public declarations, the constraints the
checker inferred for them, and their effect summaries. The interface is hashed. An importer is
re-checked only when the hash of an interface it reads changes---an \emph{interface firewall}: a
body edit that changes nothing in the interface stops at the module, and one that changes the
interface re-checks the importers. The on-disk cache key of a module covers its source and the
interfaces of its imports, so anything that is a function of those needs no new key.

Inferred interface facts churn. In a measurement of beni's inferred \code{where} constraints
(Appendix~\ref{app:method}), one class of body edit---adding an equality test on a parameter---changed
the interfaces of 50 of 90 unannotated public declarations of the core library, and of none of the
annotated ones.

The compiler's design targets are a checking throughput above 250\,000 lines per second per core,
an 800\,ms cold build of 100\,000 lines, and warm rebuilds of 15\,ms after a body edit and 60\,ms after
a signature edit. Its output must not depend on thread timing: identifiers are assigned from the
input before any parallel work starts, and a determinism test builds a corpus with one thread and
with eight, twice each, and compares every output byte.

\section{How lists are represented today}
\label{sec:bg-lists}

\code{List} is beni's only sequence type; the language deliberately has no separate array type. A
list value is one of three JavaScript forms: a plain array; a \emph{view}, an object
\code{\{ b, o, length \}} giving an offset into a base array (used for the \code{rest} of a pattern, in
$O(1)$); or a 32-way radix trie with claimable head and tail buffers. A view is always a suffix of its
base. Every write is persistent: \code{List.set} copies a plain array of up to 256 elements and
path-copies a trie above that; \code{List.push} copies short arrays (under 32 elements) and claims
the trie's tail otherwise; \code{List.pop} copies short arrays and makes a new header otherwise;
\code{insertAt} and \code{removeAt} copy; prepending through \code{[ x, …xs ]} claims a free head slot,
copies a short array, or converts. Readers outside the library use three facts---\code{length},
\code{Array.isArray} and a \code{\$plain()} method that materialises a plain array---so that the trie
and the view are invisible to them.

This design gives every program persistence at a cost that is modest per operation but paid
everywhere: the trie's code ships with every program that edits a list, every generic read path
can meet a non-plain form, and a list handed to JavaScript must first be flattened. A hand-written
JavaScript program would simply write \code{a[i] = v}.

\paragraph{Identity promises.} The current list library returns, in many cases, an input itself or a
view of it rather than a new array: \code{filter} that keeps every element, \code{map} whose every
result is \code{===} its element, \code{slice} or \code{take} of everything, \code{xs ++ []},
\code{[] ++ ys}, \code{concat} with one non-empty list, \code{reverse} and the sorts of a list
shorter than two, \code{drop xs 0}, and \code{drop} and \code{tail}, which return views. Out-of-range
\code{set}, \code{swap} and \code{removeAt}, \code{swap xs i i}, a \code{set} that writes an element
\code{===} to the one already there, and \code{pop} of the empty list return their input unchanged;
and a prepend whose element is \code{===} to the one in a claimed head slot returns the base array
itself. The renderer uses some of these promises to skip work by a pointer test. Under the rule every
one of them must either appear in the function's summary or be dropped (\S\ref{sec:model-demands}).

\paragraph{Building loops.} The backend compiles a function that builds its result in tail position
modulo a list constructor---a function whose recursive call appears inside \code{[ x, …recursive call ]}---into a loop
that pushes onto an \emph{array builder} and closes it at the exit. For example
\begin{Verbatim}
appendTo xs ys = case xs of
    [] → ys
    [ x, …rest ] → [ x, …appendTo rest ys ]
\end{Verbatim}
becomes a loop that pushes each \code{x} onto a builder and, at the exit, copies the elements of
\code{ys} onto it---or, when the builder is still empty, returns \code{ys} itself. The rewrite is
mandatory and part of the language's cost contract. Today it is mandatory for the stack's sake: an
Elm-style building recursion would otherwise recurse once per element. The builder is linear: it is
written only by the loop that made it.

\paragraph{Loops and scalar views.} A self tail call compiles to a loop that rebinds the parameters
in place. Each new argument reads the parameters as they were; arguments that make a call run in
parameter order, but the loop may rebind a parameter early when no argument still to come reads
it. When such a loop walks its list with
\code{[ x, …rest ]}, the backend replaces the view by an integer offset into the base array (a
\emph{scalar view}), so no view object is allocated per step.

\section{The release optimiser}
\label{sec:bg-release}

A build with \code{--release} runs an optimiser over the emitted code. Reachability-driven dead-code
elimination runs on every build and removes declarations no root reaches. Under \code{--release} the
optimiser also removes local bindings that are never read when their right-hand side is pure,
propagates constant arguments, inlines a private function called from one site into that site, folds
a binding that is read once into its single use, gives record fields and top-level names short names,
and prints compactly. It does not specialise a function per call site. It keeps every binding that
may be impure or may suspend, never reorders two calls, and never duplicates an argument's
evaluation. The language promises that \emph{a release build behaves exactly as the development build
does}, with no exception.

One rule matters here. The language definition lets the backend choose where an evaluation that
makes no call runs relative to another such evaluation---arithmetic, a comparison, a field read,
building a record, a literal, a name---because such evaluations change nothing and cannot fail; an
evaluation that makes any call keeps its place. The single-use fold relies on it: it folds a binding
whose initialiser is an atom or a chain of member reads into its use when every statement in between
is another such binding, and it folds a read of a beni value even into a statement whose earlier
operands make calls, on the ground that nothing can change a beni value. After lowering, two list
reads are such member reads: \code{List.length xs} is \code{xs.length}, and an element read in the
library's loops is \code{xs[i]}. So \code{n = List.length xs} followed by \code{h (List.push xs 1) n}
may be emitted as \code{h(push(xs, 1), xs.length)}, which reads the length after the push. The ground
holds only while every list is immutable; assumption A5 (\S\ref{sec:soundness-assumptions})
restates it.

\section{The browser: TEA and two back ends}
\label{sec:bg-browser}

\paragraph{The Elm Architecture.} A browser program is written in The Elm Architecture (TEA): a
\code{Model} type holding the application state, an \code{init} value (today a top-level value, not
a function), an
\code{update : Msg, Model → Model} function (or, with commands, one returning
\code{Model × Cmd Msg}) and a \code{view} from model to markup. A runtime owned by the platform
calls \code{update} for each message, replaces the current model with the result, and brings the
page up to date. \code{update} is where nearly all list edits in an application happen, and it is
where the runtime and the program both hold the model. Markup is part of the language: \code{view}
is written with JSX-like elements whose vocabulary a platform declares, with control-flow elements
such as \code{<For each=\{⋯\}>} for lists. An event handler attribute such as
\code{onClick=\{Select model.id\}} names the message to send when the event fires.

\paragraph{The template runtime.} beni's current browser back end compiles each markup site into a
cloned template and patches dynamic holes with a reference check per hole, in the style of
SolidJS's compiler. It keeps the last value of every hole, and it keeps each list it rendered so that
it can skip the next walk when the list is identical. Every rendered list is therefore held by the
runtime, invisibly to the program.

\paragraph{The direct platform.} A second, experimental back end---which we call the \emph{direct}
platform---compiles away \code{view} at run time altogether. It compiles each message handler into a
JavaScript function that runs the corresponding \code{update} arm and then writes the DOM directly. It
keeps a module-level variable holding the current model, which the arm reassigns. Its run-time state
is deliberately small:
\begin{itemize}
\item \emph{slots} hold the leaf values that the page shows, and the values of \emph{derived values}:
a \code{let} in \code{view} compiles into a derived value that handlers recompute when its inputs
may have changed; a derived value's slot may hold a list, possibly a list of the model itself;
\item per-row \emph{instance arrays} hold, for each rendered row, its DOM nodes and its \emph{item}
(the list element the row shows), used only to compare items by identity when deciding which rows
moved;
\item \emph{edit scripts} update the rows of a list after a message; each script is guarded by a tag
computed from the message and the program's own conditions, never by the identity of the list;
\item \emph{listener bodies} compute a message's payload when the event fires, reading the handler's
arguments---\code{model.id}, a row's item---at the event rather than capturing them earlier.
\end{itemize}
In a development build, after each message, a verification step re-evaluates every derived value
and expects nothing to change unless the compiler is wrong. The direct platform holds no rendered
view; its instance arrays hold row items for identity comparison, and the lists it holds are those
of the model and of derived slots. That design choice is what makes the hand-over protocol of
\S\ref{sec:hard-tea} possible.

\paragraph{Write sets.} The direct platform is guided by a whole-program \emph{write-set analysis}.
For each message it computes the \emph{write set} of the corresponding \code{update} arm: the paths
of the model that the arm may change, in a path language of record fields, tuple indices,
constructor arguments and list positions, cut at a fixed depth. A handler then runs only the DOM
updates and derived-value recomputations whose reads meet the write set. The analysis interprets the
arms abstractly, tracking for each model path whether its value is still the \emph{same} value
(identity) or may have been replaced, and it summarises callees bottom-up by Kleene iteration. Its
notion of ``changed'' is ``not the same value, not \code{===}''. It does not analyse commands (their
thunks run later and are not interpreted). Its domain records that the value at one path was moved
to another during an arm, but it has no fact for two model paths that already hold the same value
when the arm begins. It was built first as a report-only pass that printed its results and changed no
emitted byte, and only then used.

\section{What in-place writes buy: a prototype study}
\label{sec:bg-measure}

Before designing a checker we measured what in-place writes would buy, in a prototype study that
transformed beni's own compiled output by hand, exactly as a compiler pass would, and ran six array
scenarios in Node~24 (five rounds; Chrome as a spot check). Appendix~\ref{app:method} gives the
method, the machine and the limits; the study is unpublished, and its numbers are indicative. The
baseline was a persistent
\emph{adaptive} array (plain copies up to 256 elements, a trie above). One of the variants, a purely
static one, wrote in place wherever last use and per-function interface summaries proved a value
unique, with parameters that a function consumes uniquely recorded in its interface and flowing
forward to its callers. It approximates the run-time half of the checker this thesis proposes, but
it is not that checker: it had no exclusive-contents flag, no ownership classes for function values
and no entry protocol, and where it could not prove uniqueness it copied rather than rejected. Its
results relative to the baseline (medians; no ordering below rests on a ratio closer to 1 than 1.25
except where a smaller one is given):

\begin{itemize}
\item \emph{Building in a fold} (collecting by \code{push}, a histogram over 1\,000 or 100\,000
buckets, a coin-change table): 0.14--0.47 times the baseline's time, the largest gain in the study,
because every write is to an accumulator nobody else holds.
\item \emph{A grid game}, steady-state ticks of 1 or 100 cells on a board of one million cells:
0.44--0.48 times; the first tick, which must copy a board that the harness kept, was slower
(1.52 times), because a whole copy costs more than converting to a trie.
\item \emph{A TEA table} of 10\,000 rows: with the runtime keeping the rendered rows, as beni's
current runtime does, no counting or static scheme could write the rows in place---every write
found the rows shared. With a runtime that hands the model over, the writes landed in place, and
bought little that was visible (0.84--0.95 times on update and swap messages), because one in-place
element write saves about a microsecond next to a 40--60\,\textmu s render walk.
\item \emph{An undo history} that keeps every version gains nothing, correctly: nothing is unique.
\item \emph{Size}: the list runtime with the trie compressed to 1\,514 bytes (brotli); without it,
706 bytes.
\end{itemize}

The same study found that a static scheme that \emph{copies} where it cannot prove uniqueness is
harmless at a message boundary and catastrophic inside a loop: applied to a callback inside a fold,
it copied on every iteration and ran 8.2 times slower than the baseline. That observation is one
reason we propose an error rather than an optimisation: a silent copy makes the cost model depend
on facts the programmer cannot see, and in the worst case turns a linear loop quadratic. The study
also found that a model record threaded through helpers must be tracked per field: without it, an
append to one list field copied every row.

\section{The question, stated precisely}
\label{sec:bg-question}

The rule is: \emph{editing a list while anything else still holds the old version is a compile
error.} Unpacked:

\begin{itemize}
\item \emph{Editing in place} means a call of one of the core library's named list writers:
\code{set}, \code{update}, \code{push}, \code{pop}, \code{swap}, \code{insertAt}, \code{removeAt},
\code{cons} and \code{append}. List syntax---a literal, a spread such as \code{[ x, …xs ]}, and
\code{++} on lists---builds a new array and is not an edit (\S\ref{sec:model-demands}).
Chapter~\ref{sec:disc-options} asks whether in-place forms of \code{map}, \code{filter},
\code{reverse}, \code{sort} and \code{take} should join the writers.
\item \emph{Anything else} is any holder of the array, of any of the kinds listed in
\S\ref{sec:model-holders}, including the operand's own binding if it is read again after the edit.
\item \emph{Still holds} means the holder is live: some execution continuing from the edit reads
through it.
\item \emph{The old version} is the very array, not an equal list.
\item \emph{A compile error} means neither a copy nor a warning: an accepted in-place edit writes in
place, always, in development and release alike. The one escape hatch is a new library function,
\code{List.copy}, which makes intentional sharing explicit.
\end{itemize}

So the rule is \emph{uniqueness at the update site, decided by liveness}.

\paragraph{What ``observe'' means.} An in-place write is observable when some computation reads
the array's elements after the write through a reference it obtained before. Three kinds of reader
exist in beni: program code (\code{List.get}, a pattern, a fold, \code{==}---list equality walks both
lists---, \code{Debug.toString}, a markup hole that shows the list); the runtime (a renderer that keeps
the last list it rendered, or the direct platform's instance items and slots); and JavaScript (a
sibling file, or a \code{Js.from}, that kept the array). Each of these is a \emph{content} holder. An
\emph{identity} holder keeps a reference only to ask \code{===} later. An in-place write does not
change identity, so an identity holder never sees a different object; what it may do is conclude
``unchanged'' from ``same object'', which is wrong after an in-place write. beni programs cannot
compare references, so this hazard belongs to the renderer, and we discharge it by an assumption on
the platforms (A4 in \S\ref{sec:soundness-assumptions}). The direct platform was designed with the
corresponding rule: it never decides that something is unchanged from the identity of an object on a
path that may be written in place.

\paragraph{Two semantics.} The precise form of the guarantee is the one O'Connor, Linares
Ar\'evalo and Rizkallah give for uniqueness types: they work ``by statically ruling out every
program where the difference between the immutable and the mutable interpretations can be
observed'' \citep{oconnor2026separation}. Aspinall, Hofmann and Kone\v{c}n\'y prove such a theorem
for a first-order language: an in-place evaluation and a ``safe'' functional evaluation agree on every
well-typed term \citep[Thm.~5.8]{aspinall2008usage}, and Linear Haskell proves one for its mutable
arrays \citep{bernardy2018linear}. We define the \emph{value semantics} $S_v$---every
update builds a new array, as beni does today---and the \emph{update semantics} $S_u$---every
accepted update writes its operand's array---and claim that, for an accepted program, they produce
the same observations: the same printed output, the same DOM writes, the same calls into JavaScript
with the same arguments marshalled by value (Chapter~\ref{sec:soundness}).

\paragraph{What is not asked.}
\begin{itemize}
\item Records, tuples and constructors are not written in place by this checker, although the
machinery extends to them (\S\ref{sec:compiler-records}). The thesis is about lists.
\item The checker does not decide whether an in-place update is \emph{worthwhile}: every accepted
update is in place. There is no cost model and no ``copy if small''.
\item It does not replace the write-set analysis; it meets it at the direct platform's entry
protocol (\S\ref{sec:hard-tea}, \S\ref{sec:compiler-writesets}).
\end{itemize}
